\section{Held-out prediction, breakdown, and intervention}

\ifnum\ResultsReady=1 The base surrogate attains median boundary error
\BaseMedianError{}, whereas the augmented surrogate attains
\AugmentedMedianError{}, improving held-out transport by
\BridgeImprovement{}.  The mean anchor residual is \AnchorMAE{}, and the
largest registered non-additive shift is \MaxInteraction{}.  These quantities
come directly from the immutable aggregate. \else This section remains
non-numerical until every registered task and adaptive boundary seed is
complete. \fi

\resultfigure{figures/figure7_predictive_bridge.pdf}{Predicted versus observed
boundaries for the blinded variant.  Marker denotes domain and horizontal
bars are hierarchical 95\% intervals on the observed boundary.}

Prediction failure is tested rather than explained post hoc: the same holdout
is scored by base and augmented surrogates.  Critical-window interventions
allocate the expensive operation only within a fixed interval around the
predicted boundary and compare against always-on, always-off, random-window,
and matched-compute controls.  \ifnum\ResultsReady=1 The registered allocation
removes on average \MeanComputeSaving{} of uniform compute while retaining
\MeanQualityRetention{} of the always-on quality gain and exceeding a random
matched-compute window by \MeanRandomGain{}. \fi

\begin{table}[t]
\centering\small
\begin{tabular}{lrrr}
\toprule Domain & Base error & Augmented error & Compute saving \\
\midrule \PredictiveDomainRows
\bottomrule
\end{tabular}
\caption{Automatically generated held-out results.  No saturated legacy order
parameter enters this table.}
\end{table}

\resultfigure{figures/figure8_theory_breakdown.pdf}{Held-out error before and
after the registered macrovariable is added.  Error bars summarize variation
over size and secondary control, not pooled inner replicates.}
